\section{Klíčová návrhová rozhodnutí}
\label{sec:design:decisions}

Po vymezení požadavků je nutné rozhodnout, \emph{jakým způsobem} budou jednotlivé vlastnosti aplikace dosaženy. Tato sekce postupně rozebírá hlavní rozhodnutí, která ovlivnila podobu výsledné architektury. Formu nasazení aplikace, zdroje dat o~repozitáři, způsob získávání aktuálních dat společně se schopností doplnit chybějící historii a~volbu nástroje pro vizualizaci.

Další úžeji zaměřené rozhodnutí, volba konkrétního databázového systému, je probíráno v~\customref{Sekci}{sec:implementation:db_choice}.

\subsection{Forma nasazení aplikace}
\label{subsec:decisions:deployment}

Z~požadavku na webovou aplikaci s~jednoduchým deploymentem vyplývá otázka, jak má být aplikace strukturována jako celek. Zvažovány byly tři varianty.

\begin{description}
  \item[Statická jednostránková aplikace s~předgenerovaným datasetem]
    Tímto způsobem je realizován například nástroj rustc-pr-tracking \sectionref{sec:rustc_pr_tracking}. Veškerá data jsou připravena dopředu a~poskytována jako statické soubory, které prohlížeč zpracovává lokálně. Výhodou je triviální nasazení (libovolný statický hosting, žádný server) a~nulová provozní režie. Nevýhodou je omezená velikost datasetu, neboť veškerá data musí být stažena do prohlížeče, a~nemožnost provádět nad daty dynamické dotazy s~uživatelskými parametry.

  \item[Backend s~databází a~oddělený frontend]
    Klient-serverová architektura, kde backend zpřístupňuje data uložená v~databázi přes \API a~frontend slouží pouze pro jejich zobrazení. Toto řešení umožňuje uchovávat data ve vysokém objemu a~efektivně se nad nimi dotazovat pomocí \SQL. Cenou je složitější nasazení (běžící server, databáze, případná kontejnerizace) a~nutnost provozovat backend nepřetržitě.

  \item[Serverless funkce nad spravovanou databází]
    Varianta, která odstraňuje potřebu vlastního serveru. Za to je ale aplikace vázána na konkrétního cloudového poskytovatele, je omezena na jím podporované databázové systémy a~lokální vývoj je složitější.
\end{description}

\textbf{Zvolené řešení.} Aplikace je realizována jako backend nad relační databází s~odděleným webovým frontendem. Hlavním důvodem je objem dat a~potřeba dynamických analytických dotazů, které ve statické variantě nejsou realizovatelné. Závislost na konkrétním cloudovém ekosystému u~serverless varianty by zároveň zhoršila přenositelnost a~znesnadnila lokální vývoj. Jednoduchost lokálního nasazení je řešena kontejnerizací.

\subsection{Zdroje dat o~repozitáři}
\label{subsec:decisions:data_sources}

Analytické pohledy vyžadují data o~\PR, issues, jejich historii, změněných souborech a~příslušnosti přispěvatelů k~týmům. Otázkou je, odkud tato data získávat.

\begin{description}
  \item[Pouze GitHub \API]
    GitHub by byl ideálním jediným zdrojem, ale pro tento projekt nestačí. Endpoint pro soubory změněné v~rámci \PR by bylo nutné volat pro každý z~tisíců \PR zvlášť, což by přineslo další výrazné navýšení počtu požadavků a~zpomalilo synchronizaci kvůli síťové latenci~\cite{github_rest_api}. Limity počtu požadavků \subsectionref{subsec:github_api_limits} by tak byly rychle překročeny. Informace o~příslušnosti přispěvatelů k~oficiálním týmům Rust projektu GitHub \API neposkytuje vůbec.

  \item[Pouze lokální git repozitář]
    Klon repozitáře poskytuje přímý přístup ke struktuře commitů a~jejich obsahu bez síťové režie. Sám o~sobě však neobsahuje metadata o~\PR a~issues (stavy, štítky, recenze, komentáře, časovou osu událostí), bez nichž nelze požadované analytické pohledy realizovat.

  \item[Kombinace více zdrojů]
    Každý zdroj poskytne část dat: GitHub \API poskytuje metadata o~\PR a~issues včetně jejich časové osy, lokální git repozitář slouží k~výpisu souborů změněných v~merge commitech a~Rust Team \API \sectionref{sec:rust_team_api} mapuje přispěvatele do týmů.
\end{description}

\textbf{Zvolené řešení.} Aplikace kombinuje všechny tři zdroje. Lokální git odbourává tisíce \API požadavků, které by jinak mohly narážet na rate limit. Team \API doplňuje informace o týmech a jejich přispěvatelích. GitHub \API zůstává primárním zdrojem metadat o~\PR a~issues a~jejich časové ose.

\subsection{Aktuálnost dat a~self-healing}
\label{subsec:decisions:freshness}

Z~požadavků na vždy čerstvá data a~na schopnost doplnit data po výpadku vyplývá otázka, jakým mechanismem aplikace získává nové i~chybějící události z~GitHubu.

\begin{description}
  \item[Webhooky]
    GitHub při každé relevantní události zašle \HTTP požadavek na předem registrovaný endpoint~\cite{github_webhooks}. Přístup je vhodný pro reaktivní zpracování v~reálném čase, avšak má dvě podstatná omezení. Vyžaduje veřejně dostupný endpoint, což komplikuje lokální vývoj i~nasazení v~prostředích bez veřejné IP adresy. A~z~principu neposkytuje historická data události, ke kterým došlo před registrací webhooku nebo v~době výpadku aplikace, jsou nenávratně ztraceny.

  \item[Pravidelný polling REST \API stavu \PR a~issues]
    Aplikace by v~pravidelných intervalech načítala aktuální stav \PR a~issues a~detekovala změny. Implementačně jednoduché, ale nezachytí mezistavy mezi dvěma odečty (například krátkodobě přidaný a~odebraný štítek) a~neposkytuje úplnou časovou osu událostí.

  \item[Timeline events \API]
    GitHub pro každý \PR a~issue zpřístupňuje endpoint vracející úplný chronologický seznam zaznamenaných událostí~\cite{github_rest_api}. Tímto způsobem lze provést jak prvotní naplnění historických dat, tak inkrementální doplňování. Cenou je vysoký počet \API požadavků při prvotní synchronizaci.
\end{description}

\textbf{Zvolené řešení.} Aplikace používá Timeline events \API jako primární zdroj historie a~aktuálních událostí. Synchronizace je inkrementální a~opakovaně spustitelná, čímž je dosaženo požadované self-healing vlastnosti. Po libovolně dlouhém výpadku je při dalším spuštění historie automaticky doplněna. Webhooky byly zamítnuty kvůli závislosti na veřejném endpointu a~neschopnosti dohnat data zpětně. Polling není dostatečný pro rekonstrukci kompletní časové osy.

\subsection{Vizualizace dat}
\label{sec:design:visualization}

Ze specifikovaných analytických dotazů vyplývají dvě hlavní formy vizualizace: časové řady stavů \PR (dotazy Q1 a~Q2) vyžadující interaktivní spojnicové grafy a~hierarchická data o~upravených souborech (dotaz Q7) vyžadující prostorovou agregaci stromové struktury. Otázkou je, jaký nástroj pro vizualizaci zvolit.

\begin{description}
  \item[Monitorovací platforma (Grafana)]
    Grafana\footnote{\url{https://grafana.com}} je v~oblasti monitoringu zavedeným standardem s~přímou podporou PostgreSQL. Pro tento projekt se však ukázala jako nevhodná ze dvou důvodů. Integrace analytických \SQL dotazů obsahujících podmíněnou logiku a~vícenásobné klauzule \code{WITH} naráží na limity při předávání dynamických parametrů. A~vizualizace hierarchických textových dat (cesty k~souborům seskupené do adresářové struktury) není nativními panely Grafany podporována. Grafana nabízí pouze plochý tabulkový výpis, který u~stovek záznamů postrádá přehlednost.

  \item[Vlastní frontendový klient]
    Jednostránková webová aplikace komunikující s~backendem přes \REST \API. Poskytuje plnou kontrolu nad předáváním parametrů dotazů a~umožňuje implementovat libovolné vizualizace včetně treemap pro stromová data a~interaktivních grafů s~možností skrývat jednotlivé řady. Cenou je nutnost implementace a~údržby vlastního kódu.
\end{description}

\textbf{Zvolené řešení.} Vlastní frontendový klient ve formě jednostránkové webové aplikace. Hlavním důvodem je potřeba vizualizací, které hotové monitorovací nástroje nepodporují, zejména prostorová agregace souborových cest do treemap a~plná flexibilita při parametrizaci dotazů. Konkrétní pohledy a~jejich implementace jsou popsány v~\customref{Sekci}{sec:implementation:visualization}.
